\documentclass[10pt,a4paper]{amsart}
\usepackage[margin=19mm]{geometry}
\usepackage{amsmath,amssymb,mathtools}
\usepackage[scr=rsfs,cal=euler]{mathalfa}
\usepackage{xfrac}
\usepackage[unicode,hidelinks,bookmarks=false]{hyperref}
\newcommand{\ie}{i.e.\ }
\newcommand{\cf}{cf.\ }
\newcommand{\resp}{respectively\ }
\newcommand{\Subsection}{\subsection}
\providecommand{\nobreakdash}{\nobreak\hbox{-}}
\setlength{\emergencystretch}{2em}
\multlinegap0pt
\usepackage{bm}
\usepackage{amssymb}
\usepackage{xcolor}
\usepackage{algorithm}\usepackage{algpseudocode}
\newtheorem{definition}{Definition}
\newtheorem{assumption}{Assumption}
\newtheorem{theorem}{Theorem}
\title{Heaviside Low-Rank Support Matrix Machine}
\author{Xianchao Xiu, Shenghao Sun, Xinrong Li, Jiyuan Tao}
\date{Source: arXiv:2603.00491v1 (CC BY 4.0)}
\begin{document}
\maketitle

\noindent\emph{Source and licence.} Heaviside Low-Rank Support Matrix Machine. Version arXiv:2603.00491v1 (CC BY 4.0), distributed by arXiv under the Creative Commons Attribution 4.0 International licence, http://creativecommons.org/licenses/by/4.0/, as recorded in the arXiv licence metadata for this exact version and shown on the abstract page https://arxiv.org/abs/2603.00491v1. Full licence notice: \texttt{licenses/arxiv-2603.00491-CC-BY-4.0.txt}.\par\smallskip

\noindent\textit{Editorial scope.} Counted unit: the article's theorem theorem (source0.tex lines 532--533). The article's complete original proof of it is delivered with it (source0.tex lines 535--578, one \texttt{proof} environment of 44 lines). No mathematics is written, repaired or re-derived: the statement and the proof are the article's own lines, in the article's own notation, and no retained line is substituted or re-flowed.  Retained uncounted as the source-local context the proof needs: lines 450--457; lines 493--503; lines 524--526.  Nothing was omitted from any retained span, so the package carries no omission record.  No retained line contains a citation, so the wrapper carries no bibliography and no bibliography entry.  No retained line contains a figure environment, an image-inclusion command or drawing code, so no image is delivered and the package declares no asset dependency.  Editorial formatting only, in addition: an AMS \texttt{amsart} preamble that restates the article's own declarations and macros used by the retained lines, namely the environments \texttt{definition}, \texttt{assumption}, \texttt{theorem} on separate counters, together with the standard packages those lines need.  The retained environments print in this document as Definition 1, Assumption 1, Theorem 1 in order of appearance, whereas the article prints the same items with the numbering of its own full document.  The complete native source of the article as distributed in the arXiv e-print for this exact version is bundled unchanged as \texttt{sources/195493/source0.tex}.
\begin{equation}\label{RSMM3} 
	\begin{aligned}
        \min\limits_{\hat{\mathbf{W}},\mathbf{z}}\quad
		&\frac{1}{2}\langle\hat{\mathbf{W}},\hat{\mathbf{W}}\rangle+\beta\|{\mathbf{z}}_+\|_0\\
		{\rm s.t.}\quad &\mathbf{1}+\mathscr{A}(\hat{\mathbf{W}})=\mathbf{z},\\
		\quad & \text{rank}(\hat{\mathbf{W}}) \leq r.\\
	\end{aligned}
\end{equation}

\begin{definition}\label{stat}
	A point $(\hat{\mathbf{W}},\mathbf{z})$ is called a Karush–Kuhn–Tucker (KKT) point of problem \eqref{RSMM3} if there exists a Lagrangian multiplier $\bm{\lambda}$ such that
	\begin{equation*}
		\begin{cases}
			-\partial_{(\hat{\mathbf{W}},\mathbf{z})}\mathscr{L}(\hat{\mathbf{W}},\mathbf{z},\bm{\lambda})
			\in \mathrm{N}_{\mathcal{F}}(\hat{\mathbf{W}},\mathbf{z}),\\[4pt]
			\mathrm{rank}(\hat{\mathbf{W}}) \leq r,\\[4pt]
			\mathbf{1}+\mathscr{A}(\hat{\mathbf{W}})-\mathbf{z}=0.
		\end{cases}
	\end{equation*}
\end{definition}

\begin{assumption}\label{ass1}
When $s=r$, the matrices $\mathbf{T}_{\hat{\mathbf{W}}}^i$, $i=1,2,\cdots,m$, are linearly independent. When $s<r$, the matrices $\mathbf{R}_{\hat{\mathbf{W}}}^i$, $i=1,2,\cdots,m$, are linearly independent.
\end{assumption}

\begin{theorem}[Necessary optimality condition] Suppose that Assumption \ref{ass1} holds. Then a local minimizer of problem \eqref{RSMM3} is a KKT point.
\end{theorem}

\begin{proof}
	Let $\mathbf{T}=(\hat{\mathbf{W}}, \mathbf{z})$ be a local minimizer of  problem $\eqref{RSMM3}$ and denote $f(\mathbf{T})$ as the objective function. Invoking the generalized Fermat's theorem, it derives that  
	$-\partial_\mathbf{T} f(\mathbf{T})\in\mathrm{N}_{\mathcal{F}}(\mathbf{T})$. Then it has
	\begin{equation}\label{T1}
		\begin{cases}
			-\nabla_{\hat{\mathbf{W}}}f(\mathbf{T})\in\mathrm{N}_{\mathcal{L}\cap \mathcal{R}}(\hat{\mathbf{W}}),\\
			-\partial_\mathbf{z} f(\mathbf{T})\in\mathrm{N}_{\mathcal{L}}(\mathbf{z}).
		\end{cases} 
	\end{equation}
Next, consider the following two cases, which correspond to the two scenarios of Assumption \ref{ass1} respectively.
\begin{description}
    \item[(Case 1)] If $\text{rank}(\hat{\mathbf{W}})=r$, it derives that  $\mathrm{N}_{\mathcal{L}\cap \mathcal{R}}(\hat{\mathbf{W}})=\mathrm{N}_{\mathcal{L}}(\hat{\mathbf{W}})+\mathrm{N}_{\mathcal{R}}(\hat{\mathbf{W}})$.  Then,	
\begin{equation*}\label{W1}
		-\nabla_{\hat{\mathbf{W}}} f(\mathbf{T})\in\mathrm{N}_{\mathcal{L}}(\hat{\mathbf{W}})+\mathrm{N}_{\mathcal{R}}(\hat{\mathbf{W}}).
	\end{equation*}

    \item[(Case 2)]     If $\text{rank}(\hat{\mathbf{W}})<r$, it shows that $\mathrm{N}_{\mathcal{L}\cap \mathcal{R}}(\hat{\mathbf{W}})=\mathrm{N}_{\mathcal{L}}(\hat{\mathbf{W}})$. Note that
	\begin{equation*}
		\mathrm{N}_{\mathcal{L}}(\hat{\mathbf{W}})=\{\sum_{i=1}^m(\bm{\lambda})_i y_i\mathbf{X}_i \mid (\bm{\lambda})_i\in\mathbb{R}, i=1,\cdots,m \}.
	\end{equation*}
    
	Relation \eqref{T1} indicates that there  exist $\bm{\lambda}$ such that $-\nabla_{\hat{\mathbf{W}}} f(\mathbf{T})+\sum_{i=1}^m(\bm{\lambda})_i y_i\mathbf{X}_i\in\mathrm{N}_{\mathcal{R}}(\hat{\mathbf{W}})$. Then,
	\begin{equation}\label{W2}
    -\partial_{\hat{\mathbf{W}}} \mathscr{L}(\mathbf{T},\bm{\lambda})= -\nabla_{\hat{\mathbf{W}}}f(\mathbf{T})+\sum_{i=1}^m(\bm{\lambda})_i y_i\mathbf{X}_i\in\mathrm{N}_{\mathcal{R}}(\hat{\mathbf{W}}).
    \end{equation}
\end{description}
Similarly, since $\mathrm{N}_{\mathcal{L}}(\mathbf{z})=\mathbb{R}^m$, relation \eqref{T1} implies that 
    	\begin{equation}\label{N1}
        -\partial_\mathbf{z} \mathscr{L}(\mathbf{T},\bm{\lambda})\in\mathrm{N}_{\mathcal{L}}(\mathbf{z}).
            \end{equation}
Combing \eqref{W2} and \eqref{N1}, it holds that
	\begin{equation*}
		\begin{cases}
			-\partial_{\hat{\mathbf{W}}} \mathscr{L}(\mathbf{T},\bm{\lambda})\in\mathrm{N}_{\mathcal{R}}(\hat{\mathbf{W}}),\\
			-\partial_\mathbf{z} \mathscr{L}(\mathbf{T},\bm{\lambda})\in\mathrm{N}_{\mathcal{L}}(\mathbf{z}).
		\end{cases} 
	\end{equation*} 	
Therefore, it has 
	\begin{equation*}
-\partial_{({\hat{\mathbf{W}},\mathbf{z}})}\mathscr{L}(\hat{\mathbf{W}},\mathbf{z},\bm{\lambda})
	\in\text{N}_{\mathcal{F}}(\hat{\mathbf{W}},\mathbf{z}).
	\end{equation*} 
From Definition \ref{stat},  it is concluded that $\mathbf{T}$ is a KKT point. The proof is completed.
\end{proof}

\end{document}
